\hypertarget{chapitre-0-python}{%
\section{Chapitre 0 : Python}\label{chapitre-0-python}}

\textbf{Avant de commencer}

Avant de commencez à appendre Python, vous aurez besoin
d\textquotesingle outils pour écrire votre code et le compiler.

Pour à la fois gérer différentes versions de Python, installer et isoler
vos dépendances, compiler et conditionner vos modules et structurer vos
projets, installez \href{}{uv} et vérifiez que la commande suivante
affiche la version de {uv} installée et retourne \emph{exit code} de
`0`:

\begin{Shaded}
\begin{Highlighting}[]
\NormalTok{uv {-}{-}version}
\end{Highlighting}
\end{Shaded}

Étant donné un projet existant, installez tout ce qu\textquotesingle il
faut avec une seule commande:

\begin{Shaded}
\begin{Highlighting}[]
\NormalTok{uv sync}
\end{Highlighting}
\end{Shaded}

Pour gérer vos fichier et rédiger votre code, à la fois {Visual Studio
Code} (VSCode) et JetBrains \href{}{PyCharm} sont des excellents
éditeurs pour Python.

Pour simplifier l\textquotesingle apprentissage, ce guide présentera des
instructions à exécuter en ligne de commandes dans un terminal (comme
\href{}{BaSH} ou \href{}{PowerShell}). Si vous connaissez des manière
d\textquotesingle arriver aux même résultats en utilisant une interface
graphique, n\textquotesingle hésitez pas à vous en servir.

Le but de ce chapitre est que vous appreniez à coder comme un
informaticien, jonglant avec des pratiques et approches tirées des
mathématiques, du langage et de l\textquotesingle ingénierie. Nous
essaierons non seulement de développer des compétences techniques de
programmation, mais aussi des compétences de conception et de
communication de systèmes parfois complexes ou abstraits.

part*
